\section{模型评价与改进方向}
\subsection{模型特点与适用性}
问题一的主要特点是保留原生变长观测，同时区分词级、词组级与未确定关系。双顺序约束使被选词组在原词索引和声音时间上均不重叠，逐维有效时长聚合避免缺失音高对均值的稀释。该方法适用于需要保留局部缺失和对应不确定性的音视频特征整理。

问题二将缺失处理置于编码之前，并用可用比例补充融合输入，能够在统一接口下比较不同缺失模式。类别加权不增加推理参数，其收益主要体现在训练方案选择阶段的中性类F1与Macro-F1，代价是整体Accuracy及正向类F1下降。

问题三通过显式单项与交互项组织融合输出，能够在一次前向中取得可核算的模态和位置分数。可靠性系数、模态净贡献与局部读出分别记录，便于检查它们之间的差异；原始强度与类别一致性处理后的强度并列保留。

\subsection{局限与改进方向}
首先，问题一的75.98\%是时间候选覆盖率，独立词界准确性尚未测量。后续应在未参与规则开发的样本中标注词与词组边界，同时评价候选精确率、覆盖率和端点误差，而不是只追求候选数量。

其次，现有模型比较受到开发划分和选择过程限制。历史开发曾查看测试表现，问题三主要为单种子实验；问题二历史类别权重比较中固定模型的分组Bootstrap差值区间也包含零。进一步比较应按原视频分组，分开模型选择与最终评价，并补充多种子结果。问题二已补充类别加权、固定4轮下的缺失增强对照，后续需检验其跨种子和跨划分的稳定性；问题三需比较基础加性结构与可靠性缩放的稳定性。

再次，分解一致性只检验模型读出核算。BERT位置包含上下文，当前高分词元不直接等于人类情感原因；原特征索引也不等于精确媒体时间。后续可比较屏蔽高分、随机和低分位置后的重新预测结果，并结合人工证据标注评价解释质量。媒体对应待确认或定位不足的样本继续保留明确状态。

最后，问题二已同步最终模型的预测、指标、消融与核验记录，但本地尚未取得对应权重和标准化参数，独立推理复现仍需补齐模型记录所列文件。本次文稿同步依据新增实验记录完成，未在本地重新训练模型、开展人工标注或运行扰动实验。

\section{结论}
本文完成了题目要求的三类输出：附件一100条样本的多模态特征与时序关系，附件三30条情感预测，以及附件四20条预测和解释记录。

问题一通过词级筛选与词组动态规划，将候选覆盖扩展到1468/1932，并在固定候选条件下消除词组间不兼容关系。问题二最终模型的验证集Macro-F1为0.6102、MAE为0.5861，已更新附件三全部30条预测。缺失实验显示文本移除的影响较大；固定训练方案下的消融支持保留连续缺失增强，以兼顾强度预测与重度共同缺失时的分类表现。问题三给出可核算的局部融合读出，验证集最终MAE为0.5397，专项样本中18条取得可回看的关键范围。三项结果分别回答特征如何对应、缺失时如何预测和模型分数如何形成；其评价范围限定于已有观测、划分及实验记录。

\paragraph{辅助工具使用说明}
文稿整理与构建代码使用OpenAI Codex辅助，已有记录标注模型为GPT-6，使用日期为2026年9月26日至27日。本次数值核对依据既有预测和实验文件，未新增模型训练。模型版本发布日期及此前研究阶段的完整使用记录尚待据实补充。
